\documentclass[10pt,a4paper]{amsart}
\usepackage[margin=19mm]{geometry}
\usepackage{amsmath,amssymb,mathtools}
\usepackage[scr=rsfs,cal=euler]{mathalfa}
\usepackage{xfrac}
\usepackage[unicode,hidelinks,bookmarks=false]{hyperref}
\newcommand{\ie}{i.e.\ }
\newcommand{\cf}{cf.\ }
\newcommand{\resp}{respectively\ }
\newcommand{\Subsection}{\subsection}
\providecommand{\nobreakdash}{\nobreak\hbox{-}}
\setlength{\emergencystretch}{2em}
\multlinegap0pt
\usepackage{bm}
\usepackage{bbm}
\usepackage{dsfont}
\usepackage{xspace}
\usepackage{cancel}
\usepackage{mathtools}
\usepackage{amsthm}
\makeatletter
\@ifundefined{theorem}{\newtheorem{theorem}{Theorem}}{}
\makeatother
\makeatletter
\expandafter\ifx\csname proof\endcsname\relax
\newenvironment{proof}[1][Proof]{\noindent\textbf{#1.} }{\ \rule{0.5em}{0.5em}}
\fi
\makeatother
\title[Kernels for composition of positive linear operators]{Kernels for composition of positive linear operators}
\author{Ulrich Abel, Ana Maria Acu, Margareta Heilmann,, Ioan Ra\c{s}a}
\date{Source: arXiv:2506.23579v1 (CC BY 4.0)}
\begin{document}
\maketitle

\noindent\emph{Source and licence.} Kernels for composition of positive linear operators. Version arXiv:2506.23579v1 (CC BY 4.0), distributed under the Creative Commons Attribution 4.0 International licence, as recorded in the arXiv licence metadata for this exact version and shown on the abstract page https://arxiv.org/abs/2506.23579v1. Full rights notice: licenses/arxiv-2506.23579-CC-BY-4.0.txt.\par\smallskip

\noindent\textit{Editorial scope.} Counted unit: the Theorem whose source label is \texttt{None} in the article (source0.tex lines 343--354, source label \texttt{None}), delivered together with the article's own complete original proof of it (source0.tex lines 356--451, one \texttt{proof} environment of 96 lines). No mathematics is written, repaired or re-derived: statement and proof are the article's own text line for line. Required context retained uncounted: P1 (lines 237--246); eq.w2 (lines 249--255); eq.w3 (lines 249--255); eq.ww1 (lines 274--290). Every definition, display and label of those lines is retained unchanged. Omitted from this package and not redistributed: 1--236, 247--248, 256--273, 291--342, 355--355, 452--810. Nothing inside a retained excerpt is omitted or altered: every retained body is delivered exactly as the article's lines are written, so no omission receipt of any kind accompanies this package. Citations: the retained excerpts carry no citation command at all, so no bibliography entry of the article is transcribed and no reference is redistributed. Reference adaptation: none was needed and none was made; every reference inside the delivered unit resolves inside it, so no unresolved reference or citation is produced. Printed slips kept literal: no printed word, symbol, hypothesis or number of the counted statement or of its proof was altered. Layout and class: the article's own class is \texttt{article} with packages including color, hyperref, amssymb, amsmath, amsfonts, fontenc, enumerate, graphicx; the wrapper is the lane's pinned \texttt{amsart} editorial wrapper at 10pt on A4, which restates the theorem-like environments the retained text needs and adds the article's own macro definitions that the retained text needs (none), with extra wrapper packages (bm, bbm, dsfont, xspace, cancel, mathtools). The delivered numbering is the unit's own. Assets: no retained span contains a figure environment, a graphics-inclusion command, a PDF-page inclusion, a table or a TikZ picture; no image file is copied into this package, the package declares no asset dependency and no image file is redistributed with it. Licence: version arXiv:2506.23579v1 of this article is distributed under the Creative Commons Attribution 4.0 International licence, as recorded in the arXiv licence metadata for this exact version and shown on the abstract page https://arxiv.org/abs/2506.23579v1. A full rights notice is delivered at licenses/arxiv-2506.23579-CC-BY-4.0.txt. Characters: every retained excerpt is pure ASCII, so the delivered engine, XeLaTeX with its default Latin Modern text font, reports no missing character.
\begin{theorem}
\label{P1} \label{corollary-kernel-K_m,n[-1]}For $m,n\geq 0$, the kernel of
the composition of classical Bernstein--Durrmeyer operators $M_{m}\circ
M_{n} $ has the representation 
\begin{equation*}
K_{m,n}\left( x,y\right) =\frac{\left( m+1\right) !\left( n+1\right) !}{%
\left( m+n+1\right) !}\sum_{k=0}^{\min \{m,n\}}\binom{m}{k}\binom{n}{k}%
\sum_{\ell =0}^{k}p_{k,\ell }\left( x\right) p_{k,\ell }\left( y\right) .
\end{equation*}
\end{theorem}

\begin{align}
& \displaystyle \sum_{k=0}^{n}p_{n,k}(x)=1,  \label{eq.w1} \\
& \int_{0}^{1}p_{n,k}\left( t\right) dt=\frac{1}{n+1}, \text{ \quad } k=0,\ldots,n ,  \label{eq.w2} \\
& p_{m,\mu }\left( t\right) p_{n,\nu }\left( t\right) =\frac{\binom{m}{\mu }%
\binom{n}{\nu }}{\binom{m+n}{\mu +\nu }}p_{m+n,\mu +\nu }\left( t\right) .
\label{eq.w3}
\end{align}

\begin{eqnarray}
&&\sum_{\nu =0}^{n}p_{n,\nu }\left( y\right) (m+n-\mu -\nu )^{\underline{%
m-\mu }}(\mu +\nu )^{\underline{\mu }}  \label{eq.ww1} \\
&=&\left( \frac{\partial }{\partial s}\right) ^{m-\mu }\left( \frac{\partial 
}{\partial t}\right) ^{\mu }\left. s^{m-\mu }t^{\mu }\sum_{\nu
=0}^{n}p_{n,\nu }\left( y\right) s^{n-\nu }t^{\nu }\right \vert _{s=t=1} 
\notag \\
&=&\left. \left( \frac{\partial }{\partial s}\right) ^{m-\mu }\left( \frac{%
\partial }{\partial t}\right) ^{\mu }\left( s^{m-\mu }t^{\mu }\left[
ty+s\left( 1-y\right) \right] ^{n}\right) \right \vert _{s=t=1}  \notag \\
&=&\sum_{i=0}^{\mu }\binom{\mu }{i}\sum_{j=0}^{m-\mu }\binom{m-\mu }{j}\frac{%
\mu !}{i!}\frac{\left( m-\mu \right) !}{j!}
n^{\underline{i+j}}   %% Change 25.06.2025
y^{i}\left( 1-y\right) ^{j}  \notag \\
&=&\displaystyle \sum_{i=0}^{\mu }{\binom{\mu }{i}}\sum_{j=0}^{m-\mu }{%
\binom{m-\mu }{j}}{\binom{n}{i+j}}\mu !(m-\mu )!p_{i+j,i}(y).  \notag
\end{eqnarray}%

\begin{theorem}
For $r=3$ the kernel of the composition of classical Bernstein--Durrmeyer $%
M_{n_3}\circ M_{n_2}\circ M_{n_1}$ operators has the representation 
\begin{align*}
K_{n_3,n_2,n_1} \left( x,y\right) &=\frac{\left( n_3+1\right) !\left(
n_2+1\right) !\left( n_1+1\right) !\left( n_3+n_2+n_1+1\right) !}{\left(
n_3+n_2+1\right) !\left( n_3+n_1+1\right) !\left( n_2+n_1+1\right) !} \\
&\times \sum_{k= 0}^{\min \{n_1,n_2,n_3\}}\frac{\binom{n_3}{k}\binom{n_2}{k}%
\binom{n_1}{k}}{\binom{n_3+n_2+n_1+1}{k}}\sum_{\ell =0}^{k}p_{k,\ell }\left(
x\right) p_{k,\ell }\left( y\right).
\end{align*}
\end{theorem}

\begin{proof}
By Theorem \ref{P1}, we have 
\begin{equation*}
K_{n_{2},n_{1}}(t,y)=\dfrac{(n_{2}+1)!(n_{1}+1)!}{(n_{2}+n_{1}+1)!}%
\displaystyle \sum_{k\geq 0}{\binom{n_{2}}{k}}{\binom{n_{1}}{k}}\sum_{\ell
=0}^{k}p_{k,\ell }(t)p_{k,\ell }(y).
\end{equation*}%
Using the recurrence formula for the kernel and Theorem \ref{P1} we get 
\begin{align*}
K_{n_{3},n_{2},n_{1}}(x,y)& =\displaystyle%
\int_{0}^{1}K_{n_{2},n_{1}}(t,y)K_{n_{3}}(x,t)dt \\
& =\displaystyle \int_{0}^{1}\dfrac{(n_{2}+1)!(n_{1}+1)!}{(n_{2}+n_{1}+1)!}%
\sum_{k\geq 0}{\binom{n_{2}}{k}}{\binom{n_{1}}{k}}\sum_{\ell
=0}^{k}p_{k,\ell }(t)p_{k,\ell }(y) \\
& \times (n_{3}+1)\sum_{\mu =0}^{n_{3}}p_{n_{3},\mu }(x)p_{n_{3},\mu }(t)dt
\\
& =\dfrac{(n_{2}+1)!(n_{1}+1)!}{(n_{2}+n_{1}+1)!}\sum_{k\geq 0}{\binom{n_{2}%
}{k}}{\binom{n_{1}}{k}}\sum_{\ell =0}^{k}p_{k,\ell }(y) \\
& \times (n_{3}+1)\sum_{\mu =0}^{\mu }p_{n_{3},\mu }(x)\int_{0}^{1}p_{k,\ell
}(t)p_{n_{3},\mu }(t)dt.
\end{align*}%
Using (\ref{eq.w2}) and (\ref{eq.w3}) we obtain 
\begin{eqnarray*}
\lefteqn{K_{n_{3},n_{2},n_{1}}(x,y)} \\
&=&\dfrac{(n_{3}+1)!(n_{2}+1)!(n_{1}+1)!}{(n_{2}+n_{1}+1)!}\sum_{k\geq 0}{%
\binom{n_{2}}{k}}{\binom{n_{1}}{k}}\sum_{\ell =0}^{k}p_{k,\ell }(y){\binom{k%
}{\ell }}\dfrac{1}{(k+n_{3}+1)!}\cdot S_{1},
\end{eqnarray*}%
where 
\begin{equation*}
S_{1}:=\displaystyle \sum_{\mu =0}^{n_{3}}p_{n_{3},\mu }(x)(k+n_{3}-\ell
-\mu )^{\underline{k-l}}(\ell +\mu )^{\underline{\ell }}.
\end{equation*}

In the same way as in (\ref{eq.ww1}) we derive 
\begin{equation*}
S_{1}=\displaystyle \sum_{i=0}^{\ell }{\binom{\ell }{i}}\sum_{j=0}^{k-\ell }{%
\binom{k-\ell }{j}}{\binom{n_{3}}{i+j}}\ell !(k-\ell )!p_{i+j,i}\left(
x\right) .
\end{equation*}%
Thus, by interchanging the order of summation we get 
\begin{align*}
K_{n_{3},n_{2},n_{1}}(x,y)& =\dfrac{(n_{3}+1)!(n_{2}+1)!(n_{1}+1)!}{%
(n_{2}+n_{1}+1)!} \\
& \times \sum_{k\geq 0}{\binom{n_{2}}{k}}{\binom{n_{1}}{k}}\dfrac{k!}{%
(n_{3}+k+1)!}\sum_{i=0}^{k}\dfrac{1}{i!}\sum_{j=0}^{k-i}\dfrac{1}{j!}{\binom{%
n_{3}}{i+j}}p_{i+j,i}(x)S_{2},
\end{align*}%
where 
\begin{align*}
S_{2}& :=\displaystyle \sum_{\ell =i}^{k-j}p_{k,\ell }(y)\ell ^{\underline{i}%
}(k-\ell )^{\underline{j}}=\dfrac{k!}{(k-j-i)!}y^{i}(1-y)^{j} \\
& ={\binom{k}{i+j}}i!j!p_{i+j,i}(y).
\end{align*}%
Now 
\begin{align*}
K_{n_{3},n_{2},n_{1}}(x,y)& =\dfrac{(n_{3}+1)!(n_{2}+1)!(n_{1}+1)!}{%
(n_{2}+n_{1}+1)!}\sum_{k\geq 0}{\binom{n_{2}}{k}}{\binom{n_{1}}{k}}\dfrac{k!%
}{(n_{3}+k+1)!} \\
& \times \sum_{i=0}^{k}\sum_{j=0}^{k-i}{\binom{n_{3}}{i+j}}p_{i+j,i}(x){%
\binom{k}{i+j}}p_{i+j,i}(y).
\end{align*}%
Collecting all terms with $i+j=\ell $ and interchanging the order of
summation leads to 
\begin{align*}
K_{n_{3},n_{2},n_{1}}(x,y)& =\dfrac{(n_{3}+1)!(n_{2}+1)!(n_{1}+1)!}{%
(n_{2}+n_{1}+1)!}\sum_{k\geq 0}{\binom{n_{2}}{k}}{\binom{n_{1}}{k}}\dfrac{k!%
}{(n_{3}+k+1)!} \\
& \times \sum_{\ell =0}^{k}{\binom{n_{3}}{\ell }}{\binom{k}{\ell }}\sum_{\nu
=0}^{\ell }p_{\ell ,\nu }(x)p_{\ell ,\nu }(y) \\
& =\dfrac{(n_{3}+1)!(n_{2}+1)!(n_{1}+1)!}{(n_{2}+n_{1}+1)!}\sum_{\ell \geq 0}%
{\binom{n_{3}}{\ell }}\sum_{\nu =0}^{\ell }p_{\ell ,\nu }(x)p_{\ell ,\nu }(y)
\\
& \times \sum_{k\geq \ell }{\binom{n_{2}}{k}}{\binom{n_{1}}{k}}\dfrac{k!}{%
(n_{3}+k+1)!}{\binom{k}{\ell }}.
\end{align*}%
For the inner sum we get

\begin{eqnarray*}
&&\sum_{k\geq \ell }{\binom{n_{2}}{k}}{\binom{n_{1}}{k}}\dfrac{k!}{%
(k+n_{3}+1)!}{\binom{k}{\ell }} \\
&=&\dfrac{n_{2}!n_{1}!}{\ell !}\dfrac{1}{(n_{2}+n_{3}+1)!(n_{1}-\ell )!}%
\sum_{k}{\binom{n_{1}-\ell }{k}}{\binom{n_{2}+n_{3}+1}{n_{2}-k-\ell }} \\
&=&\dfrac{n_{2}!n_{1}!}{\ell !(n_{1}-\ell )!(n_{2}+n_{3}+1)!}{\binom{%
n_{1}+n_{2}+n_{3}-\ell +1}{n_{2}-\ell }}.
\end{eqnarray*}%
Therefore, 
\begin{align*}
K_{n_{3},n_{2},n_{1}}(x,y)& =\dfrac{(n_{3}+1)!(n_{2}+1)!(n_{1}+1)!}{%
(n_{3}+n_{2}+1)!(n_{3}+n_{1}+1)!(n_{2}+n_{1}+1)!} \\
& \times \sum_{\ell \geq 0}(n_{1}+n_{2}+n_{3}-\ell +1)!\ell !{\binom{n_{3}}{%
\ell }}{\binom{n_{2}}{\ell }}{\binom{n_{1}}{\ell }}\sum_{\nu =0}^{\ell
}p_{\ell ,\nu }(x)p_{\ell ,\nu }(y)
\end{align*}%
and the proof is complete.
\end{proof}

\end{document}
